% then I would not build one giant Java/Python application. I would build an AI platform with Java as the core enterprise/backend layer and Python as the specialized AI/agent layer.

The architecture could look like this:

                         ┌──────────────────────┐
                         │   Web / Mobile UI    │
                         │ React / Next.js etc. │
                         └──────────┬───────────┘
                                    │
                              HTTPS / WebSocket
                                    │
                         ┌──────────▼───────────┐
                         │      API Gateway     │
                         │       Java           │
                         └──────────┬───────────┘
                                    │
              ┌─────────────────────┼─────────────────────┐
              │                     │                     │
       ┌──────▼──────┐       ┌──────▼──────┐       ┌─────▼─────┐
       │ Auth / User │       │   Chat      │       │   SaaS    │
       │   Service   │       │   Service   │       │  Service  │
       │    Java     │       │    Java     │       │   Java    │
       └─────────────┘       └──────┬──────┘       └───────────┘
                                    │
                              ┌─────▼──────┐
                              │ AI Gateway │
                              │    Java    │
                              └─────┬──────┘
                                    │
                    ┌───────────────┼────────────────┐
                    │               │                │
             ┌──────▼──────┐ ┌──────▼──────┐ ┌──────▼──────┐
             │ OpenAI      │ │ Python AI   │ │ RAG Engine  │
             │ API         │ │ Services    │ │             │
             └─────────────┘ └──────┬──────┘ └──────┬──────┘
                                    │                │
                             ┌──────▼────────────────▼─────┐
                             │       AI / Agent Layer      │
                             │ Python + OpenAI             │
                             └──────────────┬──────────────┘
                                            │
                    ┌───────────────────────┼────────────────────┐
                    │                       │                    │
              ┌─────▼─────┐          ┌─────▼─────┐       ┌─────▼─────┐
              │ PostgreSQL│          │ Vector DB  │       │   Redis   │
              │           │          │            │       │           │
              └───────────┘          └────────────┘       └───────────┘


# 1. Java should be your foundation

% I'd make Java/Spring Boot the main application platform.

% Something like:

% ## ai-platform/
% │
% ├── gateway/
% ├── auth/
% ├── users/
% ├── organizations/
% ├── billing/
% ├── subscriptions/
% ├── chat/
% ├── conversations/
% ├── documents/
% ├── knowledge/
% ├── ai/
% ├── agents/
% ├── coding/
% ├── audit/
% ├── notifications/
% └── admin/


# 2. Python becomes your AI laboratory

Python is extremely useful for the parts where the AI ecosystem is strongest.

For example:

python-ai/
│
├── agents/
│   ├── coding_agent/
│   ├── research_agent/
│   ├── support_agent/
│   └── data_agent/
│
├── rag/
│   ├── ingestion/
│   ├── chunking/
│   ├── embeddings/
│   ├── retrieval/
│   └── reranking/
│
├── evaluation/
├── document_processing/
├── tool_execution/
└── shared/

Your Java backend can communicate with these services through:

Java
  │
  │ REST / gRPC
  ▼
Python AI Service

You could also use a message broker for long-running jobs:

Java
 │
 ▼
Kafka / RabbitMQ
 │
 ▼
Python Worker


3. Your chatbot

Your chatbot could be:

User
 │
 ▼
Java API
 │
 ├── authenticate user
 ├── check subscription
 ├── load conversation
 ├── check permissions
 └── send AI request
          │
          ▼
      AI Gateway
          │
          ▼
       OpenAI
          │
          ▼
       Response
          │
          ▼
      PostgreSQL

The important thing is that Java owns the business rules.

For example:

if (!subscriptionService.canUseAI(user)) {
    throw new UsageLimitException();
}

if (!authorizationService.canAccessConversation(user, conversation)) {
    throw new ForbiddenException();
}

Then the AI layer handles the actual model interaction.
4. Your RAG system

This is where Python becomes particularly useful.

Suppose a company uploads:

employee_handbook.pdf
security_policy.pdf
engineering_documentation.pdf
contracts.pdf

Your pipeline could be:

                 Upload
                   │
                   ▼
              Java API
                   │
                   ▼
              File Storage
                   │
                   ▼
             Python Worker
                   │
        ┌──────────┼──────────┐
        ▼          ▼          ▼
     Parse       Chunk      Metadata
        │          │          │
        └──────────┼──────────┘
                   ▼
              Embeddings
                   │
                   ▼
              Vector DB

Then when the user asks:

    "What is our vacation policy?"

Question
   │
   ▼
Java
   │
   ▼
RAG Service
   │
   ├── query embedding
   ├── vector search
   ├── metadata filtering
   ├── reranking
   └── context construction
   │
   ▼
OpenAI
   │
   ▼
Answer + citations

This gives you a serious enterprise RAG architecture.
5. Coding assistant

You can make a separate service:

Coding Assistant
       │
       ├── Repository ingestion
       ├── Code indexing
       ├── Symbol extraction
       ├── Semantic search
       ├── Context selection
       ├── Code generation
       ├── Code review
       └── Test generation

For example:

User:
"Why is authentication failing?"

             │
             ▼

       Coding Agent
             │
       ┌─────┼─────┐
       ▼     ▼     ▼
    Search  Git   Tests
     Code   Repo  / Build
       │     │     │
       └─────┼─────┘
             ▼
           OpenAI
             │
             ▼
        Explanation

Eventually you could support:

    "Fix this bug."

The agent could inspect the repository, identify relevant files, propose a patch, run tests in a sandbox, inspect failures, and iterate.

That execution environment must be heavily sandboxed. Don't let an LLM directly execute arbitrary commands on your production server.
6. AI agents

This is where I'd use Python heavily.

You could have:

                    Agent Manager
                         │
        ┌────────────────┼────────────────┐
        │                │                │
        ▼                ▼                ▼
  Research Agent   Coding Agent     Support Agent
        │                │                │
        ▼                ▼                ▼
     Tools            Tools            Tools

For example:

Research Agent
 ├── web_search
 ├── database_search
 ├── document_search
 ├── calculator
 └── report_generator

An agent shouldn't have unrestricted access to your system.

Instead, expose controlled tools:

tools = [
    search_documents,
    get_customer,
    create_ticket,
    query_database,
]

The model requests a tool.

Your application decides whether it is allowed.

That's a very important security boundary.

OpenAI's Python Agents SDK is specifically designed around agents, tools, sessions, human-in-the-loop workflows, handoffs, guardrails and tracing.
7. Your SaaS platform

Now we get to the interesting part.

Don't design this as:

User → AI

Design it as:

                    SaaS Platform
                         │
          ┌──────────────┼──────────────┐
          ▼              ▼              ▼
       Tenant A       Tenant B       Tenant C
          │              │              │
       Users          Users          Users
          │              │              │
       Projects       Projects       Projects
          │              │              │
       Documents      Documents      Documents
          │              │              │
          └──────────────┼──────────────┘
                         ▼
                    AI Platform

Every important entity should have tenant ownership.

For example:

Organization
    │
    ├── Users
    ├── Projects
    ├── Conversations
    ├── Documents
    ├── Agents
    ├── API Keys
    ├── Usage
    └── Billing

Then you can sell:

Free
 └── 10,000 AI tokens/month

Pro
 └── 1,000,000 tokens/month

Business
 └── 10,000,000 tokens/month

Enterprise
 ├── SSO
 ├── RBAC
 ├── Audit logs
 ├── Private knowledge bases
 ├── Custom retention
 └── Advanced security

8. Enterprise security

This is where I agree with your instinct to use Java — but not because Java magically makes an application secure.

You want security architecture like:

                    API Gateway
                         │
                         ▼
                    JWT/OAuth2
                         │
                         ▼
                  Authentication
                         │
                         ▼
                  Authorization
                         │
               ┌─────────┴─────────┐
               ▼                   ▼
             RBAC              Tenant Check
               │                   │
               └─────────┬─────────┘
                         ▼
                    AI Service
                         │
                         ▼
                    Tool Policy
                         │
                         ▼
                   OpenAI / Tools

I'd also include:

    Secrets manager
    Encryption at rest
    TLS everywhere
    Audit logging
    Rate limiting
    Tenant isolation
    RBAC
    API key rotation
    Input validation
    Output validation
    Prompt-injection defenses
    Tool permission policies
    Sandboxed code execution
    Dependency scanning
    Security monitoring

9. Don't let Python bypass Java security

This is a very important architectural rule.

Don't do:

Python
  │
  ├── database
  ├── users
  ├── billing
  ├── permissions
  └── OpenAI

Instead:

                 Java
                  │
          Security boundary
                  │
        ┌─────────┴─────────┐
        ▼                   ▼
    PostgreSQL          Python AI
                            │
                            ▼
                         OpenAI

Python should receive only what it needs.

For example:

{
  "tenantId": "...",
  "userId": "...",
  "requestId": "...",
  "task": "RAG_QUERY",
  "query": "What is our vacation policy?",
  "allowedTools": [
    "search_knowledge_base"
  ]
}

Not:

{
  "databasePassword": "...",
  "adminToken": "..."
}

Obviously.
10. Database architecture

I'd start with:

PostgreSQL
     │
     ├── users
     ├── organizations
     ├── memberships
     ├── projects
     ├── conversations
     ├── messages
     ├── documents
     ├── agents
     ├── tools
     ├── usage
     ├── subscriptions
     └── audit_logs

Then add a vector database depending on your requirements.

You might use PostgreSQL with a vector extension initially rather than immediately introducing another distributed system.

Redis can handle things such as:

Redis
 ├── caching
 ├── sessions
 ├── rate limits
 ├── temporary state
 └── queues

11. Your AI Gateway

This is probably the most important service I'd create.

Instead of every service talking directly to OpenAI:

Chat ──────────────┐
RAG ───────────────┤
Coding ────────────┤
Agents ────────────┤
                  ▼
             AI Gateway
                  │
                  ▼
                OpenAI

The AI Gateway can handle:

AI Gateway
│
├── authentication
├── authorization
├── tenant limits
├── model selection
├── token usage
├── cost tracking
├── retries
├── timeouts
├── streaming
├── structured output
├── tool policies
├── logging
├── tracing
└── OpenAI API

That gives you one place to control your entire AI infrastructure.
12. Java package design

I'd make the Java side something like:

com.yourcompany.aiplatform

├── auth
│   ├── controller
│   ├── service
│   ├── security
│   └── repository
│
├── tenant
│
├── user
│
├── chat
│   ├── controller
│   ├── service
│   ├── domain
│   └── repository
│
├── ai
│   ├── gateway
│   ├── model
│   ├── provider
│   ├── streaming
│   ├── tools
│   └── usage
│
├── rag
│   ├── document
│   ├── ingestion
│   ├── retrieval
│   └── citation
│
├── agent
│   ├── controller
│   ├── service
│   ├── policy
│   └── execution
│
├── coding
│
├── billing
│
├── audit
│
└── common

13. Python package design

And:

ai_engine/

├── agents/
│   ├── base.py
│   ├── coding.py
│   ├── research.py
│   └── support.py
│
├── rag/
│   ├── ingestion.py
│   ├── chunking.py
│   ├── embeddings.py
│   ├── retrieval.py
│   └── reranking.py
│
├── tools/
│   ├── search.py
│   ├── database.py
│   ├── github.py
│   └── documents.py
│
├── safety/
│   ├── guardrails.py
│   └── policies.py
│
├── evaluation/
│
├── api/
│
└── main.py

14. How Java and Python communicate

For the first version, keep it simple:

Java ──HTTPS──> Python

Use REST APIs.

Later, for high-performance internal services:

Java ◄──── gRPC ────► Python

For asynchronous processing:

Java
 │
 ▼
Kafka
 │
 ▼
Python Worker

Don't introduce Kafka, Kubernetes, service meshes, 20 microservices, etc. on day one.

That creates a huge amount of complexity before you have a product.
15. The technology stack I'd choose
Backend

Java 21/25+
Spring Boot
Spring Security
Hibernate/JPA
PostgreSQL
Redis

AI

OpenAI API
OpenAI Java SDK
Python
OpenAI Python SDK / Agents SDK

The official Python SDK exposes the Responses API as well, while the Agents SDK provides a higher-level agent runtime.
RAG

PostgreSQL + vector capability
Object storage
Embedding model
Retrieval
Reranking

Infrastructure

Docker
Kubernetes later
Nginx/API Gateway
Cloud object storage
Secrets Manager
Monitoring
OpenTelemetry

Frontend

React
Next.js
TypeScript

16. Most importantly: don't start with everything

You said "everything", but I would build it in stages.
Phase 1 — AI foundation

Build:

Java Spring Boot
       │
       ▼
AI Gateway
       │
       ▼
OpenAI

Features:

    Authentication
    Chat
    Streaming
    Conversation history
    Usage tracking

Phase 2 — SaaS

Add:

Users
Organizations
Projects
RBAC
Subscriptions
Usage limits
Billing

Phase 3 — RAG

Add:

Document upload
       ↓
Document processing
       ↓
Chunking
       ↓
Embeddings
       ↓
Vector search
       ↓
Context
       ↓
OpenAI

Phase 4 — Agents

Add:

Agent
 ├── tools
 ├── memory
 ├── planning
 ├── execution
 ├── permissions
 └── human approval

Phase 5 — Coding assistant

Add:

Git repository
      ↓
Code indexing
      ↓
Code search
      ↓
Coding agent
      ↓
Patch
      ↓
Sandbox
      ↓
Tests
      ↓
Review

Phase 6 — Enterprise

Then add:

SSO
SCIM
RBAC
Audit logs
Data retention
Tenant isolation
Enterprise policies
Security monitoring
Compliance controls

The big architectural idea

I wouldn't think of your project as:

    "Java + Python + OpenAI."

I'd think of it as:

    A multi-tenant AI platform, with Java controlling the application/security/business layer and Python providing specialized AI/agent capabilities.

Something like:

                       YOUR AI PLATFORM
                              │
       ┌──────────────────────┼──────────────────────┐
       │                      │                      │
     SaaS                  Enterprise              AI
       │                      │                      │
       ▼                      ▼                      ▼
   Users/Billing          Security/RBAC        AI Gateway
   Projects               Audit/SSO             │
   Teams                   Policies              │
                                                ▼
                              ┌─────────────────┼────────────────┐
                              │                 │                │
                           Chatbot             RAG             Agents
                              │                 │                │
                              └─────────────────┼────────────────┘
                                                │
                                         Coding Assistant
                                                │
                                                ▼
                                           OpenAI API

And Java doesn't have to compete with Python. Let each language do what it is good at.

For the OpenAI integration specifically, I'd put a clean AIProvider / AIService abstraction in Java, so your application isn't tightly coupled to one model provider. The official Java SDK already supports the Responses API, streaming, structured outputs and function calling, which makes Java quite capable of being the main AI backend.

If you're serious about building this as one large project, the next useful step is to design the actual production-grade project from zero: repository structure → Spring Boot services → Python AI service → PostgreSQL schema → RAG pipeline → agent architecture → authentication/RBAC → Docker → deployment.

